\documentclass[10pt, twocolumn]{article}
\usepackage{graphicx}
\usepackage{amsmath}
\usepackage{amssymb}
\usepackage{amsthm}
\usepackage{mathtools}

\theoremstyle{plain}
\newtheorem{theorem}{Theorem}[section]
\newtheorem{lemma}[theorem]{Lemma}
\newtheorem{corollary}[theorem]{Corollary}
\newtheorem{proposition}[theorem]{Proposition}

\theoremstyle{definition}
\newtheorem{definition}[theorem]{Definition}
\newtheorem{example}[theorem]{Example}

\theoremstyle{remark}
\newtheorem{remark}[theorem]{Remark}

\title{%
  \vspace{-1.5em}%
  \large\textbf{On Theorems for Hartley's Law}%
}

\author{A. E. Mahrouss\\amlal@nekernel.org}
\date{\today}

\begin{document}

\twocolumn[{%
  \maketitle
  \vspace{-0.5em}
  \vspace{0.8em}
  \rule{\linewidth}{0.4pt}
  \vspace{1em}
}]

\section{Equations \& Theorems}

\begin{theorem}
Let:
$\forall S \in \mathbb{R}; \quad Baud(MIN) < S < Baud(MAX)$:
\begin{align*}
	Freq(S) &= \frac{Hi(S)}{Lo(S)} Weight(x)\\
	-\frac{Freq(S)}{Weight(S)} &= \frac{Hi(S)}{Lo(S)}
\end{align*}
By Hartley's Law, the weight of $S$ over its frequency $F$, may be equal of the division between the highest signal and lowest signal of $S$.
\end{theorem}

\nocite{*}

\bibliographystyle{acm}
\bibliography{Bell_System_Hartley}

\end{document}